\section{Kisah}

Perkembangan web membuat informasi dapat disebarkan dengan sangat cepat.
Halaman web, media sosial, forum, dan situs promosi dapat memuat berbagai
jenis konten hanya dalam hitungan detik. Di sisi lain, kemudahan penyebaran
informasi juga membuka ruang bagi munculnya konten berisiko, salah satunya
konten yang berkaitan dengan judi daring atau \textit{judi online} (judol).
Konten seperti ini sering kali muncul dalam bentuk kata kunci eksplisit,
variasi penulisan, angka yang disisipkan ke dalam kata, hingga teks yang
tertanam di dalam gambar.

Deteksi manual terhadap konten judol tidak selalu praktis. Sebuah halaman web
dapat memiliki banyak elemen teks, tautan, gambar promosi, dan istilah yang
sengaja dimodifikasi agar tidak mudah ditemukan oleh pencarian biasa. Sebagai
contoh, kata seperti \texttt{gacor} dapat ditulis menjadi \texttt{g4cor},
\texttt{slot} dapat muncul sebagai bagian dari frasa promosi, dan nama situs
dapat ditulis dengan pola huruf diikuti angka seperti \texttt{ZEUS222}.
Variasi semacam ini membuat pencocokan kata secara eksak saja tidak cukup.

Permasalahan tersebut dapat dipandang sebagai persoalan pencocokan string.
Sebuah sistem perlu membaca teks pada halaman web, membandingkannya dengan
daftar kata kunci yang dicurigai, mengenali pola tertentu menggunakan
\textit{regular expression}, serta mendeteksi bentuk penulisan yang mirip
melalui pendekatan \textit{approximate string matching}. Selain itu, karena
sebagian konten promosi muncul dalam gambar, sistem juga perlu memiliki
kemampuan \textit{Optical Character Recognition} (OCR) agar teks pada gambar
dapat diproses oleh pipeline deteksi yang sama.

Pada Tugas Besar ini, dikembangkan sebuah ekstensi Chromium bernama
\textbf{Judol Detector}. Ekstensi ini melakukan pemindaian terhadap teks yang
terlihat pada halaman web, mendeteksi indikasi konten judol menggunakan
kombinasi algoritma pencocokan string, menandai teks yang terdeteksi secara
langsung pada halaman, serta menampilkan statistik hasil deteksi melalui popup
ekstensi. Program juga menyediakan mode blur untuk menyamarkan konten yang
terdeteksi dan fitur OCR untuk memindai teks yang tertanam pada gambar.

\section{Latar Belakang}

Pencocokan string merupakan salah satu persoalan mendasar dalam ilmu komputer.
Dalam bentuk paling sederhana, pencocokan string bertujuan menemukan kemunculan
suatu pola dalam teks. Persoalan ini banyak digunakan dalam mesin pencari,
editor teks, bioinformatika, keamanan siber, dan sistem moderasi konten.
Algoritma klasik seperti Knuth-Morris-Pratt (KMP) dan Boyer-Moore dirancang
untuk melakukan pencocokan eksak secara efisien dengan mengurangi jumlah
perbandingan karakter yang tidak perlu \cite{cormen2009introduction}.

Pada konteks deteksi konten judol, pencocokan eksak berguna untuk menemukan
kata atau frasa yang sudah diketahui, seperti \texttt{slot}, \texttt{rtp},
\texttt{bonus deposit}, atau \texttt{situs judi}. Namun, konten promosi sering
menggunakan variasi penulisan untuk menghindari deteksi sederhana. Penulisan
seperti \texttt{H0KI88}, \texttt{G4COR}, atau \texttt{SLOT777} tetap dapat
dimaknai oleh manusia sebagai istilah judol, meskipun tidak sama persis dengan
kata kunci dasar. Oleh karena itu, diperlukan pendekatan lain seperti
\textit{regular expression} dan \textit{fuzzy matching}.

\textit{Regular expression} dapat digunakan untuk mendeteksi pola tekstual
yang terstruktur, misalnya gabungan huruf dan angka yang sering digunakan pada
nama situs atau kode promosi \cite{friedl2006regex}. Sementara itu,
\textit{Levenshtein Distance} mengukur jarak edit antara dua string berdasarkan
jumlah operasi penyisipan, penghapusan, dan penggantian karakter yang
diperlukan untuk mengubah satu string menjadi string lain \cite{levenshtein1966}.
Dengan memberikan bobot khusus pada karakter yang mirip secara visual, sistem
dapat mengenali variasi seperti huruf \texttt{o} yang diganti angka
\texttt{0}, huruf \texttt{a} yang diganti angka \texttt{4}, atau huruf
\texttt{s} yang diganti angka \texttt{5}.

Web modern juga tidak hanya memuat teks dalam bentuk node DOM. Banyak konten
promosi disisipkan dalam gambar, banner, atau elemen visual. Untuk menangani
kasus tersebut, Judol Detector menggunakan OCR berbasis Tesseract.js untuk
mengambil teks dari gambar yang terlihat pada halaman \cite{tesseractjs}.
Teks hasil OCR kemudian diproses menggunakan pipeline deteksi yang sama dengan
teks DOM.

Pemilihan bentuk ekstensi browser memungkinkan deteksi dilakukan langsung pada
halaman yang sedang dibuka pengguna. Dengan Manifest V3, ekstensi Chromium
dapat menjalankan \textit{content script} untuk berinteraksi dengan halaman,
serta menyediakan popup sebagai antarmuka ringkas untuk menampilkan status dan
statistik pemindaian \cite{chrome_extensions}. Pendekatan ini membuat sistem
tidak perlu mengunduh halaman melalui server terpisah, melainkan bekerja di
sisi klien pada halaman yang sedang aktif.

\section{Deskripsi Tugas}

Program Judol Detector adalah ekstensi Chromium untuk mendeteksi indikasi
konten judi daring pada halaman web. Program memindai teks yang terlihat pada
halaman, menjalankan algoritma pencocokan terhadap daftar kata kunci, menandai
hasil deteksi pada halaman, dan menampilkan ringkasan statistik melalui popup.

Secara umum, alur kerja program adalah sebagai berikut:

\begin{enumerate}
    \item Pengguna membuka halaman web yang ingin diperiksa.
    \item Pengguna membuka popup ekstensi dan menjalankan proses pemindaian.
    \item \textit{Content script} mengumpulkan node teks yang terlihat pada DOM.
    \item Program menjalankan pipeline deteksi, yaitu pencocokan pola dengan
    RegEx dan pencocokan fuzzy dengan Weighted Levenshtein.
    \item Jika OCR diaktifkan, program juga memindai gambar yang terlihat,
    melakukan preprocessing citra, dan menjalankan pipeline deteksi terhadap
    teks hasil OCR.
    \item Teks atau gambar yang terdeteksi diberi penanda visual pada halaman.
    \item Popup menampilkan jumlah hasil deteksi, waktu eksekusi, perbandingan
    algoritma, kata kunci paling sering muncul, dan ringkasan OCR.
\end{enumerate}

Program dikembangkan menggunakan TypeScript, React, Vite, Tailwind CSS, dan
Tesseract.js. TypeScript digunakan untuk menjaga struktur data deteksi,
React digunakan untuk membangun antarmuka popup, Vite digunakan sebagai build
tool, Tailwind CSS digunakan untuk styling, dan Tesseract.js digunakan untuk
fitur OCR.

\subsection{Masukan}

Masukan yang digunakan oleh program terdiri atas beberapa komponen berikut:

\begin{enumerate}
    \item \textbf{Halaman web aktif}, yaitu halaman yang sedang dibuka pada
    browser Chromium dan menjadi target pemindaian.

    \item \textbf{Teks DOM yang terlihat}, yaitu kumpulan node teks yang dapat
    dibaca dari struktur DOM halaman dan tidak berada pada elemen yang harus
    diabaikan seperti \texttt{script}, \texttt{style}, \texttt{textarea}, dan
    \texttt{input}.

    \item \textbf{Daftar kata kunci}, yaitu kumpulan istilah yang berkaitan
    dengan judol, seperti \texttt{slot}, \texttt{gacor}, \texttt{rtp},
    \texttt{jackpot}, \texttt{casino online}, dan berbagai variasi frasa
    promosi lain.

    \item \textbf{Gambar halaman}, yaitu elemen gambar yang terlihat pada
    viewport dan dapat diproses melalui OCR setelah melalui tahap preprocessing.

    \item \textbf{Pengaturan popup}, yaitu pilihan pengguna seperti menjalankan
    ulang pemindaian, mengaktifkan OCR, dan mengaktifkan blur terhadap konten
    yang terdeteksi.
\end{enumerate}

\subsection{Keluaran}

Keluaran yang dihasilkan program adalah sebagai berikut:

\begin{enumerate}
    \item \textbf{Highlight pada halaman}, yaitu penandaan visual pada teks
    yang terdeteksi sebagai indikasi konten judol.

    \item \textbf{Blur konten}, yaitu penyamaran visual terhadap teks atau
    gambar yang terdeteksi ketika fitur blur diaktifkan atau ketika gambar
    hasil OCR dinyatakan positif.

    \item \textbf{Tooltip deteksi}, yaitu informasi tambahan yang muncul saat
    pengguna mengarahkan kursor ke hasil deteksi, berisi kata kunci,
    algoritma, jumlah kemunculan, dan waktu eksekusi.

    \item \textbf{Ringkasan statistik}, yaitu jumlah total deteksi, waktu
    eksekusi, jumlah deteksi per algoritma, dan daftar kata kunci yang paling
    sering ditemukan.

    \item \textbf{Ringkasan OCR}, yaitu jumlah gambar yang dipindai, jumlah
    gambar yang terdeteksi positif, jumlah gambar yang dilewati, dan waktu
    eksekusi OCR.

    \item \textbf{Status pemindaian}, yaitu status seperti siap, sedang
    memindai, terdeteksi, bersih, atau gagal memindai.
\end{enumerate}
